% !TeX root = ../papers/nc-p3.tex
\section{純関係を許す反対階数九の曲線分岐の排除}
\label{sec:rank-nine-branch-exclusion}
\begin{lemma}\label{lem:rank-nine-ternary-center-conic}
反対階数9の曲面射影のない標準正則窓で、中央支持3の順中心 $p$ の列の
第二曲線は必ず平面二次曲線である。その平面の式 $b$ を末文字とする純関係がある。
位置を逆にすれば、逆中心の第二曲線の平面の式は純関係の始文字である。
\end{lemma}
\begin{proof}
列を $p,x_1,x_2,x_3,\ldots$ とする。二直線列の下界11により $x_3$ は動く。
$x_2$ の次数は $3-f_1$ で、次数1なら次の次数が $-f_2$ となるから2または3である。
次数2では滑らかな平面二次、次数3では最初の直線上に順側欠損がなく、
$3\mathcal O(-1)\subset4\mathcal O$ の商から完全三次系となる。
その結果 $x_1,x_2$ の一般中央支持はともに4である
（定理\ref{thm:branch-binary-center-pairing}）。反対切断の一余次元の対合により
$w(p,-,-,x_3(t))$ は階数1以下である。一般に零ではない。実際 $w(p)$ の末尾は
三次元支持なのでこれを消す $x_3$ はただ一点である。

$w(p)$ を順に $S\otimes B\otimes C$ の $2\times3\times3$ テンソルとし、
末尾の縮約による三空間を $J\subset S\otimes B$ とする。$S$ 方向での
三行三列鉛筆の非零パラメータに階数1はない。もし値が $b\otimes c$ であれば、
他方の値を $B/kb$ に射影して消す $C^*$ の非零元を取れる。これで末尾を縮約すると
$S\otimes kb$ に非零の元を得る（末尾支持は全 $C$）。二つの連続純関係の中央文字が
$b$ となるので不可能である。階数0は先頭の完全支持に反する。

$J$ 中の上の動く階数1行列は、射影的に一定でないとしてもその $B$ 因子は一定である。
実際 $S$ 因子を固定して $B$ 因子が動けば、$S$ 因子を消す鉛筆の値は
独立な二つの $C^*$ 方向を消すので先の階数下界に反する。両因子が動く場合、
$B$ 因子が二空間を張れば、両因子の積座標は少なくとも三次元を張るため
$J$ 全体がその二空間に支持され矛盾する。$B$ 因子が三空間を張れば積座標は
四次元以上で次元に反する。ここで関数空間の積の次元下界は、一つの滑らかな点で
各空間の基底を異なる位数に取り、二つの位数集合の和集合（対ごとの和）を数えればよい。
従って $w(p,-,-,x_3(t))=a(t)\otimes b$ と書ける。
点対 $x_2(t),x_3(t)$ の条件から $b(x_2(t))=0$ で、動く第二像は非退化三次ではなく
二次曲線に一致する。一般の $t$ を固定すれば主張の純関係が得られる。
\end{proof}

\begin{theorem}\label{thm:rank-nine-branch-exclusion}
反対階数9の標準正則窓で全射影像が高々一次元なら、どの辺のどちらの射影にも
正次元のファイバーはない。従って反対階数9以下の全窓で、全次数の正則性と
添字保持同型が既出の有限判定で決まる。
\end{theorem}
\begin{proof}[二次元中心の一つ前の位置]
順中心をその中央支持によって二次元または三次元と呼ぶ。
系\ref{cor:rank-nine-nonlinear-corridors}によりこの二種類で尽きる。
二次元中心 $q$ の次位置・反対位置に、後続直線 $L_q$ と中央支持の消去直線 $N_q$ を取る。
後者は定理\ref{thm:branch-binary-center-pairing}の対をなす逆中心の前点直線である。
二項対称二乗の三関係により、$L_q$ に始まらない次の辺の点対は $N_q$ に終わる。
同じ位置の別の順中心の後続直線は $L_q$ と交わらず、その第二曲線は $N_q$ に入る。
これは二直線列の排除に反する。よって二次元中心がある位置の順中心はその一点だけである。

一つ前の位置に順中心 $p$ があれば $q\in L_p$ である。$p$ が三次元なら前補題の
$f_1=1$ と $q$ の位置での一意性から従う。$p$ が二次元なら、$q\notin L_p$ のとき
$L_q=N_p$ が必要になる。しかし $N_p$ の各点はその対をなす逆中心一つのみを後続とするので
この直線を初回とする列は動く位置一つで終わり矛盾する。
従って $p$ の第二曲線は二次曲線である（$f_1>0$ かつ二直線列はない）。
この曲線の一般点は $L_q$ の外にあり、第三曲線は $N_q$ となる。次数公式より
$f_2=3\cdot2-3-1=2$ で、二次曲線上のある順中心 $r$ の後続直線が $N_q$ と交わる。

そのような交点はない。$r$ が二次元なら $q$ の中央二空間の全文字に終わる純関係と
$r$ の順二空間の全文字に始まる純関係が連続するため、両二空間は横断し
$L_r\cap N_q=\varnothing$ である。$r$ が三次元なら、$L_r$ の一般点の中央支持は4であり、
順中心でない全点でも4である。実際直線上の表示 $w(x)=\sum K_j(x)\otimes A_j$ の
三列はそこで独立で、固定テンソルの中央支持を保つ。一方 $N_q$ のどの点も後続はただ一つで
順中心ではなく、反対の $q$ に消されるので中央支持は3以下である。
これで二次元順中心の一つ前の位置には順中心がないことが示された。
\end{proof}

\begin{proof}[三次元中心の排除]
三次元順中心が一つあれば前補題より次位置にも順中心があり、前段よりこれは三次元である。
四位置を一周すれば全位置に三次元順中心があり、二次元中心の位置での一意性から
二次元順中心は存在しない。順逆の二次元中心の対より逆中心も全て三次元である。
逆中心は少なくとも一つ存在し、逆の向きにも同じ二次曲線の議論で全四位置に存在する。

同じ位置0にある順中心 $p$ と逆中心 $z$ を取る。$p$ の第二曲線 $C_2$ の平面の式 $b$
は $R_1$ の純関係の末文字である。$z$ の三次縮約を逆の順で
$S'\otimes B'\otimes C'$（位置3,2,1）に置く。$S'$ の二次元鉛筆が正則鉛筆なら、
その全非零値は階数2以上である。前点直線 $L_z$ の外の位置3の点でのこの鉛筆は
一般には可逆で、そのとき対応する位置2の点は $B'$ の消去点のみである。
特異なパラメータは有限通りで、そこでも位置2の点はそれぞれ一本の固定直線に入る。
この制約は辺 $2\to3$ に含まれる $z$ の三縮約関係から従う。
従って $C_2$ からの第三曲線は $L_z$ に一致する。
その直線は $z$ に終わる列の末端であり、逆にたどれば $C_2$ は $z$ の第二曲線である。
逆向きの前補題より $R_2$ に同じ文字 $b$ に始まる純関係が生じる。

鉛筆が恒等的に特異なら、両支持が3で定数左右核がなく、一般階数は2である。
Kronecker分解の左右の特異ブロックの最小指数はともに正であり、
三行三列に収まるのは各一つの指数1だけである。従って逆向きの基底で
\[
                \begin{pmatrix}u&v&0\\0&0&u\\0&0&v\end{pmatrix}
\]
とできる。末尾での縮約の三関係には二次元の $S'\otimes k b'$ が入り、
元の向きではいずれも文字 $b'$ に始まる $R_2$ の純関係である。
第三曲線が $L_z$ の外を動けば、これらから $C_2$ は $b'=0$ の平面に入り
$b'$ は $b$ に比例する。$L_z$ 自身なら先と同じ逆向きの前補題を使う。
いずれも同じ $b$ を介する連続純関係になり、三次元中心は排除される。
\end{proof}

\begin{proof}[二次元中心だけの場合]
全順中心が二次元なら全逆中心も二次元である。実際逆に三次元中心が一つあれば
前段を逆に適用して矛盾する。順中心 $p$ を位置0に置き、一つ前の位置だけでなく
次位置にも順中心がない（次の中心が二次元ならその一つ前に $p$ がある）。
従って $d_2=3$ で非退化三次となる。位置2に順中心がなければ $d_3=6$ である。
平面成分の次数は3以下なので非平面的な非直線が二本連続し、
系\ref{cor:rank-nine-nonlinear-corridors}に反する。
よって順中心は位置0,2にだけ存在し、逆中心は二次元中心の対により位置1,3にだけ存在する。
位置2の順中心を $a$ とする。第三曲線は $L_a$ ではない。一般前点が動くためである。
第四位置には逆中心がなく、従って第四曲線は動いて $N_a$ に一致し、次がその終中心となる。
逆向きでも終中心の隣の位相には逆中心がないから、逆向きの第二曲線、すなわち元の第三曲線は
非退化三次である。第二曲線と連続して同じ矛盾を与える。

分岐が消えたので、系\ref{cor:all-windows-through-eight}の被覆に
系\ref{cor:all-one-dimensional-point-graphs}と定理\ref{thm:finite-rank-nine-all-marks}を適用して
階数9全体の有限判定を得る。
\end{proof}
